\documentclass[12pt,a4paper]{article}

\usepackage[a4paper, left=2.5cm, right=2.5cm, top=3cm, bottom=3cm]{geometry}  
\usepackage[brazil]{babel}  
\usepackage{float}  
\usepackage{csquotes}
\usepackage{setspace}       
\usepackage{graphicx}       
\usepackage{amsmath}        
\usepackage{hyperref}       
\usepackage{cleveref} 
\usepackage{lipsum}         
\usepackage{textcomp}       
\usepackage{booktabs}  
\usepackage{array}     
\usepackage{amssymb}
\usepackage{bookmark} 
\usepackage[utf8]{inputenc}
\usepackage[date=year,
     backend=biber,
     backref=page,
     style=numeric-comp,
     bibstyle=ieee, 
     url=false,
     isbn=true,10.1109/TPWRS.2009.2037006
     sorting=nty,
     sortcites=true]{biblatex}

\usepackage{url}
\usepackage{xurl}
	 
\doublespacing

 \AtEveryBibitem{\clearfield{pages}}
 \AtEveryBibitem{\clearfield{notes}}

\addbibresource{ic.bib}

\title{\textbf{Programação Linear Inteira}}
\author{Sofia Valverde Villas Bôas}
\date{}

\begin{document}

\maketitle

\section{Programação Linear}

\emph{Programação Linear} (PL) é uma técnica utilizada no processo de modelagem e resolução de problemas de otimização \cite{Chvatal1983}. Estes problemas envolvem a busca pela melhor solução, segundo um critério de custo ou valor, em um conjunto enumerável de possibilidades. Um programa linear na forma padrão é escrito como:

\begin{alignat}{4}
&&\text{maximize}   \quad && \sum_{j = 1}^{n} c_j x_j \label{equa:obj1}\\
&&\text{sujeito a} \quad && \sum_{j = 1}^{n} a_{ij} x_j & \leq b_i     \quad && i = 1, \ldots, m \label{res:1.1} \\
&&                        && x_j                       & \ge 0        \quad &&  j = 1, \ldots, n \label{res:2.1}
\end{alignat}

A \autoref{equa:obj1} descreve a \textit{função objetivo}, isto é, a função linear que será maximizada (ou minimizada) em um programa linear, no qual $x_j$ são chamadas de \textit{variáveis de decisão}. A \autoref{res:1.1} representa as \textit{restrições} que delimitam os possíveis valores destas variáveis. Por fim, a \autoref{res:2.1} indica uma \textit{restrição de não negatividade}, já que em muitos problemas não faz sentido atribuir valores negativos para as variáveis de decisão. 

Denomina-se \textit{solução viável} uma atribuição de valores a $x_j$ que satisfaz todas as restrições impostas ao problema. Já uma \textit{solução ótima} é uma solução viável que maximiza a função objetivo. Vale destacar que podem existir mais de uma solução ótima para o mesmo problema ou até mesmo não existir solução ótima. No caso de não existir nenhuma solução viável, o problema é chamado de \textit{inviável} e, para o caso em que as soluções viáveis permitem que a função objetivo cresça indefinidamente, os problemas são chamados de \textit{ilimitados} \cite{Chvatal1983}.

Entre as técnicas utilizadas para resolver problemas de PL, está o \textit{método simplex} \cite{Chvatal1983}. A partir de uma solução viável inicial, o \textit{simplex} realiza iterativamente o processo de pivotamento, isto é, a troca de variáveis, visando melhorar o valor da função objetivo. Apesar de ter complexidade exponencial no pior caso, o \textit{simplex} apresenta desempenho eficiente, na prática.

\section{Programação Linear Inteira}

Um \emph{Programa Linear Inteiro} (PLI) é um PL com a restrição adicional de que todas as variáveis devem assumir valores inteiros. Problemas de PLI são NP-Completo, mas, na prática, é possível resolver instâncias específicas de PLI de forma eficiente, graças ao desenvolvimento de algoritmos exatos, técnicas de otimização e heurísticas eficazes \cite{Wolsey2021}.


\subsection{Formulações}
\subsection{Otimalidade, Relaxação e Limitantes}
\subsection{Problemas Bem Resolvidos}
\subsection{Branch and Bound}
\subsection{Algoritmos de Corte}

Em muitos problemas, obter o $\mathrm{conv}(X)$ é difícil, sendo assim técnicas que obtem aproximações para o $\mathrm{conv}(X)$ se tornam importantes na resolução eficiente de instâncias dadas de determinado problema.

Nesse contexto, somos introduzidos ao conceito de \emph{inequações válidas}.

 


\begingroup
\onehalfspacing
\printbibliography
\endgroup

\end{document}
